\SourceLeaf{383}%
Bedeutung, weil sie uns einen unbestreitbaren Maßstab für die Genauigkeit der Arbeiten selbst geben. Abgesehen von Rechenfehlern können sie bis auf unmerkliche Kleinigkeiten nichts anderes sein als die zusammengefassten Fehler aus den astronomischen Längenbestimmungen, den Azimutbestimmungen und den Winkelmessungen im Dreiecksnetz. Diese Einflüsse lassen sich freilich nicht voneinander trennen, doch ist das Vorliegen des Gesamtfehlers unabhängig von den Unregelmäßigkeiten der Erdgestalt eine unleugbare Tatsache, auch wenn man seine bisherige Einschätzung der bei allen drei Arbeiten erreichten \textit{absoluten} Genauigkeit merklich herabsetzen muss. Vermutlich hat jede ihren Anteil beigetragen. obwohl wir dazu neigen, die größere Hälfte den gemessenen Dreieckswinkeln zuzuschreiben. Die meisten Arbeiten, die in diese Vergleiche eingehen, finden wir zwar in diesem Werk. Die Winkelmessungen zwischen Mailand und Superga, die in früheren Jahren von französischen Ingenieuren ausgeführt wurden, sind aber nur in abgekürzter und bereits ausgeglichener Form wiedergegeben, sodass sich der Grad ihrer Genauigkeit überhaupt nicht beurteilen lässt. In der \textit{Connaissance des tems} von 1829, S.\,288, finden wir immerhin, dass dabei Fehler in den Winkelsummen bis zu \(6''8\) vorkommen. Auch das für die Übertragung der Orientierung bei der Verbindung der Sternwarten von Mailand und Turin sehr wichtige Dreieck aus Superga und der alten und der neuen Sternwarte in Turin (S.\,254) scheint nicht ganz mit der erforderlichen Genauigkeit gemessen worden zu sein. Zur Klärung dieser so wichtigen Frage wäre es sehr wünschenswert, dass die dortigen Astronomen wenigstens zwischen den beiden Sternwarten eine neue geodätische Verbindung herstellen. Mit zwei Heliotropen und dem von beiden Sternwarten sichtbaren Standort S.\,Bernardo di Fenera wäre das äußerst leicht. Wenn von dort auch nur einer der übrigen früheren Zwischenpunkte sichtbar ist, woran kaum zu zweifeln ist, würde die ganze Arbeit nur die Messung von vier Winkeln erfordern.

Im achten Abschnitt finden wir eine sehr interessante und verdienstvolle Arbeit, eine vollständige Wiederholung der von \textsc{Beccaria} 1762\,.\,.\,.\,1764 durchgeführten Breitengradmessung. Bekanntlich ließ sich ihr Ergebnis mit den in anderen Ländern gemessenen Graden überhaupt nicht in Einklang bringen. Die Krümmung des Bogens zwischen den Endpunkten Mondovi und Andrate
